% !TeX root = ../papers/nc-p3.tex
\section{有限点窓の安定性を二つの内在的な旗だけで判定する}
\label{sec:finite-intrinsic-stability-test}
\begin{theorem}\label{thm:finite-intrinsic-stability-test}
有限Frobenius窓の四辺が有限であるとする。三次元接点があれば厳密半安定である。
それ以外では次の有限検査で安定性が決まる。長さ9以上で埋込み次元2の各局所点 $p$ に対し、
四位置で $p_i$ とその射影接平面を取る。このような点は高々二つである。
\begin{enumerate}
\item $p_i$ で非零な文字 $z_i$、接平面を切る文字 $h_i$、その間の二文字 $U_i$ を
順に重み $0,1,2$ の各部分に対応させる。
$K_i=R_i\cap(\ker p_i\otimes\ker p_{i+1})$ を、
$(\ker p_i/kh_i)\otimes(\ker p_{i+1}/kh_{i+1})$ へ写す。
全四辺で像が零なら厳密半安定である。
\item 点の長さが10以上で、局所環の極小二変数表示の二次初期式空間が
一つの一次式の平方で張られるとする。その接方向直線を各接平面内に戻し、
点、接方向直線、接平面の完全旗を取る。対応する四文字の重みを $0,1,2,3$ とする。
窓の重み6未満の全係数が零なら厳密半安定である。
\end{enumerate}
すべての点で両検査が不成立なら安定である。重み検査は旗に適合する任意の基底で行える。
したがって曲線状の点だけを持つ有限窓は、重複度によらず安定である。
三次元接点も長さ9以上の平面局所点もない場合にも同じ結論となる。
\end{theorem}
\begin{proof}[二重の中央重みと二次関係の検査]
点対の周囲で両 $z$ を1にする。グラフの接空間は両接平面の間の接線型同型のグラフなので、
六方程式の微分の階数は4である。両方で最低文字を用いない関係は正確に $K_i$ で、次元2となる。
$w(p_{i-2},\eta_{i-1})$ を、$\eta_{i-1}$ が点の代表を法として二つの接方向を動くときに取れば、
その像は $K_i$ である。実際、$(p_{i-2},p_{i-1})$ でのグラフ上の微分により
$w(p_{i-2},\eta_{i-1})=-w(\eta_{i-2},p_{i-1})$ である。
一方は末尾が、他方は先頭が点で消える関係に入る。点での評価が階数3で、
その零化点は $p_{i-1}$ だけだから接方向二つの像は独立である。

重み4未満の四語を調べる。最低文字二個が隣り合えば点対条件から零である。
反対位置に最低文字が二個なら、その間の残る少なくとも一文字は中央文字である。
対応する係数は上の $K_i$ のどちらかの端に最低文字を入れたものなので零となる。
最低文字が一個なら残る三文字は全て中央文字である。この全係数の消失は
$K_i$ の両中央文字への射影の消失とちょうど同値である。
これで第一検査が重み4未満の消失と同値であることが分かる。
第二検査と同様に数値判定の等号を与える。各検査は重み以上の部分空間のみによるので、
旗の分裂の選択には依存しない。
\end{proof}

\begin{proof}[一般の等号方向が与える旗の復元]
三次元接点があれば前定理の証明と命題\ref{prop:maximal-tangent-closed-limits}によって
正次元の点グラフを持つ閉窓へ落ち、安定ではない。
残る厳密半安定窓に、定理\ref{thm:all-finite-windows-semistable}のスカラーでない等号方向を取る。
その最低点 $p$ で、上二関係には一次部がなく、最初の四関係の一次部は独立である。
よって埋込み次元は2で、上二関係がちょうど $K_i$ である。
中央二関係の一次部は両最高文字 $h_i,h_{i+1}$ の独立な二式である。
実際、単純文字の場合にその重みは $2d+s,d+s+g$ 以上で、
最低文字を含む一次部に現れる文字は最高文字だけである。
中央二重の場合にも同じである。これで $h_i$ が実際の接平面の式となる。

中央二重の場合、$K_i$ は重み3以上であり、両中央文字への射影は零である。
単純四重み $0,d,d+s,d+s+g$ の場合も、上二関係の重みは $2d+s+g$ 以上であり、
両中央文字のどの積よりも真に大きければ第一検査が成立する。
この不等式が破れるのは $s=g$、従って $s=d=g$ の場合だけである。
ここで上二関係に残り得る中央積は $22$ のみである。
ある辺でこの係数が非零なら、全辺で非零である。
実際その重み4の関係線と反対の重み2の線の完全対合から
$2220,2202$ の係数はともに非零である（相手の $20,02$ はともに非零）。
次の切断で隣接する $22$ を読めば両隣の辺にも伝播する。

各位置の文字を $z,u,v,h$ とする。局所二変数を $u,v$ として最初の四式で他の四変数を形式的に消去する。
一次では $h,h'$ は零となり、重み2の関係の $02,20$ の両係数が非零なので
$v'$ は $v$ の非零倍となる。残る二式の二次部は一方が $vv'$ の非零倍、他方は零である。
したがって二次初期式空間は $v^2$ で張られる。この平方が決める接方向は重み
$0,1,2,3$ の元の旗の点と接平面の間の直線である。
元の等号条件は正確に第二検査の重み6未満の消失となる。
\end{proof}

\begin{proof}[点の長さ]
第一検査なら形式的な二変数への消去後、$K_i$ のどの項も一次の中央文字と二次以上の
$h,h'$ の積か最高文字どうしの積である。従って二式はいずれも位数3以上で、局所交点数は9以上となる。
第二の旗を与えた等号方向では、アフィン六変数の重みは両側で $1,2,3$、
六関係の最低重みは $1,2,3,3,4,5$ である。各変数をその重み乗で置き換えると、
形式局所環の拡大は階数36の有限自由拡大になる。置換後の各式の通常の位数は
対応する最低重み以上である。正則局所環の局所Bézout不等式で、
拡大後の完全交叉長は六つの位数の積以上、従って元の長さは $360/36=10$ 以上となる。
全長20と合わせて、検査する点数と判定の必要十分性が従う。
\end{proof}
